\documentclass[11pt]{article}
\usepackage[margin=1in]{geometry}
\usepackage{amsmath,amssymb}
\usepackage{eurosym}

\title{FRE-GY 9743 Assignment 6}
\author{Syonaa Mehra}
\date{}

\begin{document}
\maketitle

\section*{Question}

EUR/USD FX forward collateralised in GBP.

\textbf{Setup.} At settlement $T_s$ the long side receives $N^{\text{\euro}}$ EUR and pays
$N^{\text{\euro}} K$ USD, where $K$ is in USD per EUR. $X^{\$/\text{\euro}}(t)$ is the spot
rate (USD per EUR). The collateral is in GBP and earns $r_c^{\pounds}$ (SONIA).

Notation (same as the FX currency chain notes): $P^{d,col}(t,T)$ is the zero-coupon bond for a
currency $d$ cash flow collateralised in currency $col$. So $P^{\$,\$} = df^{SOFR}$,
$P^{\text{\euro},\text{\euro}} = df^{EuroSTR}$, $P^{\pounds,\pounds} = df^{SONIA}$,
$P^{\pounds,\$} = df^{[GBP,USD]}$ and $P^{\text{\euro},\$} = df^{[EUR,USD]}$.
I also write $\lambda^{\pounds} = r^{\pounds} - r^{\pounds}_c$ for the GBP funding spread.

\textbf{Which formula to use.} Pick the reporting currency $p = $ USD. The collateral currency
is GBP, so $p \neq col$ and we use the foreign collateral formula
(Eq.\ (4) in the foreign collateralisation notes):
\[
V(t) = E^{Q^d}_t\left[ e^{-\int_t^T r^d(s)ds}\, e^{\int_t^T (r^f(s) - r^f_c(s))ds}\, V(T) \right]
\]
with $d = $ USD, $f = $ GBP.

\subsection*{1. Valuation formula}

Plugging in the forward payoff:
\[
V_{\$}^{FX}(t) = N^{\text{\euro}}\, E^{Q^\$}_t\left[ e^{-\int_t^{T_s} r^\$(s)ds}\,
e^{\int_t^{T_s} \lambda^{\pounds}(s)ds} \left( X^{\$/\text{\euro}}(T_s) - K \right) \right].
\]
Split into the two terms.

\emph{Strike term.} $K$ is a constant, so it is just a USD cash flow collateralised in GBP:
\[
K\, E^{Q^\$}_t\left[ e^{-\int_t^{T_s} r^\$ ds}\, e^{\int_t^{T_s} \lambda^{\pounds} ds} \right]
= K\, P^{\$,\pounds}(t,T_s).
\]

\emph{FX term.} Change numeraire from $Q^\$$ to $Q^{\text{\euro}}$ with
\[
Z = \frac{X^{\$/\text{\euro}}(T_s)\, B^{\text{\euro}}(t,T_s)}{X^{\$/\text{\euro}}(t)\, B^{\$}(t,T_s)},
\]
where $B$ is the money market account. The $X(T_s)$ and the $e^{-\int r^\$}$ cancel and we get
\[
E^{Q^\$}_t\left[ e^{-\int_t^{T_s} r^\$ ds}\, e^{\int_t^{T_s} \lambda^{\pounds} ds}\, X^{\$/\text{\euro}}(T_s) \right]
= X^{\$/\text{\euro}}(t)\, E^{Q^{\text{\euro}}}_t\left[ e^{-\int_t^{T_s} (r^{\text{\euro}} - r^{\pounds} + r^{\pounds}_c) ds} \right]
= X^{\$/\text{\euro}}(t)\, P^{\text{\euro},\pounds}(t,T_s).
\]
This is the EUR cash flow collateralised in GBP, converted at spot.

So the value is
\[
V_{\$}^{FX}(t) = N^{\text{\euro}} \left( X^{\$/\text{\euro}}(t)\, P^{\text{\euro},\pounds}(t,T_s) - K\, P^{\$,\pounds}(t,T_s) \right).
\]
Here I am ignoring the convexity adjustment from the notes, i.e.\ assuming zero correlation
between the funding spread $\lambda$ and the payoff, so that the $P^{\cdot,\pounds}$ can be
treated as plain discount factors.

\subsection*{2. Fair forward strike}

Set $V = 0$:
\[
K = X^{\$/\text{\euro}}(t)\, \frac{P^{\text{\euro},\pounds}(t,T_s)}{P^{\$,\pounds}(t,T_s)}.
\]
This matches the general rule $FX^{f/d}(t;T_s) = X^{f/d}(t)\, P^{d,col}/P^{f,col}$.

If FX forwards are independent of the collateral currency, this is the same as the
USD-collateralised forward:
\[
K = X^{\$/\text{\euro}}(t)\, \frac{P^{\text{\euro},\$}(t,T_s)}{P^{\$,\$}(t,T_s)}.
\]
So GBP collateral does not change the fair strike. It only changes the discounting of the PV,
since $V = N^{\text{\euro}} P^{\$,\pounds}(t,T_s)\,(K^* - K)$.

\subsection*{3. Are the given discount factors enough?}

From part 1 we need $P^{\text{\euro},\pounds}$ and $P^{\$,\pounds}$. Neither is given directly,
and there is no liquid EUR/GBP cross currency market anyway. But we can use the
EUR--USD--GBP currency chain (Eq.\ (7) in the notes): assuming FX forwards don't depend on the
collateral currency,
\[
P^{\text{\euro},\pounds} = \frac{P^{\text{\euro},\$}}{P^{\pounds,\$}}\, P^{\pounds,\pounds}
= \frac{df^{[EUR,USD]}\; df^{SONIA}}{df^{[GBP,USD]}}.
\]
Doing the same thing for USD (with $\$$ in place of \euro):
\[
P^{\$,\pounds} = \frac{P^{\$,\$}}{P^{\pounds,\$}}\, P^{\pounds,\pounds}
= \frac{df^{SOFR}\; df^{SONIA}}{df^{[GBP,USD]}}.
\]
Putting these in:
\[
V_{\$}^{FX}(t) = N^{\text{\euro}}\, \frac{df^{SONIA}}{df^{[GBP,USD]}}
\left( X^{\$/\text{\euro}}(t)\, df^{[EUR,USD]} - K\, df^{SOFR} \right),
\qquad
K = X^{\$/\text{\euro}}(t)\, \frac{df^{[EUR,USD]}}{df^{SOFR}}.
\]

So yes, they are enough, with a few comments:
\begin{itemize}
  \item We still need the EUR/USD spot $X^{\$/\text{\euro}}(t)$, which is not a discount factor.
  \item The chain assumes FX forwards are independent of collateral currency, and we dropped the
        convexity adjustment (zero correlation between the xccy funding spread and the payoff).
        If that correlation isn't small we would need a model, not just curves.
  \item $df^{EuroSTR}$ is not actually used, since $df^{[EUR,USD]}$ is already given.
  \item GBP only enters through $df^{SONIA}/df^{[GBP,USD]}$, which changes the PV but not $K$.
\end{itemize}

\end{document}
